% ===== latex-recomposed:document =====
% flm: style file flmrecomposed.6b6a41b1.sty
\flmCellsTable{cc}{X[-1]X[-1]}{hspan=minimal,
  cell{1}{1}={}{m, font={\flmCellsHeaderFont}},
  cell{1}{2}={}{m, font={\flmCellsHeaderFont}},
  hline{2}={1}{.4pt,solid},
  hline{2}={2}{.4pt,solid}}{\toprule
Term&Meaning\\
Cell&A table of cells is a grid of rows and columns in which each cell occupies a rectangle of grid positions, one position or several when the cell spans rows or columns, and its content is laid out as blocks.\\
\bottomrule
}
%%
% ===== latex-recomposed:style flmrecomposed.6b6a41b1.sty =====
\NeedsTeXFormat{LaTeX2e}
\ProvidesPackage{flmrecomposed.6b6a41b1}
% --- packages ---
% flm.latex.pkg.tabularray
\RequirePackage{tabularray}
% flm.latex.pkg.xcolor
\RequirePackage{xcolor}
% --- commands and environments the library defines for the document ---
% flm.latex.def.cells
% Tables of cells. A table is written
% \flmCellsTable{<natural>}{<fitted>}{<options>}{<rows>}: <natural> and
% <fitted> are two column specifications of the tblr environment of the
% tabularray package (<natural> has c for each column, <fitted> has X[-1]
% for each column),
% <options> are the table's other tblr options (the cells' spans, alignment
% and header font, and the lines), and <rows> are the rows, the cells
% separated by & and each row ended by \\.
% By default the table is set at its natural width with <natural>; where it
% is wider than \flmCellsMaxWidth, it is set again at that width with
% <fitted>. It is set between \flmCellsBeginCenter and \flmCellsEndCenter,
% except where LaTeX sets a single line, as in a cell of another table set
% at its natural width, where no paragraph can begin. A header cell is set
% in \flmCellsHeaderFont. Redefine \flmCellsTable to set tables another way,
% or one of the four other commands to change one part of it.
% A cell's custom style <name> is written \flmCellsSet{<name>} at the start
% of the cell's text: a tabularray table command that applies the tabularray
% cell keys held in \flmCellsKeys<name>. \flmCellsStyleKeys{<name>}{<keys>}
% sets those keys unless keys are already set for <name>: define
% \flmCellsKeys<name> before loading this file, or redefine it after, to
% change how a style looks. The shipped styles green, blue, yellow and red
% set the background colors flmCellsColor<name>, which you can define before
% loading this file.
% tabularray's counter library restores LaTeX's counters after the passes in
% which tabularray sets a cell only to measure it.
% The macros with @ in their names are internals, not a customization point.
\UseTblrLibrary{booktabs}
\UseTblrLibrary{counter}
\providecommand\flmCellsHeaderFont{\bfseries}
\providecommand\flmCellsBeginCenter{\begin{center}}
\providecommand\flmCellsEndCenter{\end{center}}
\providecommand\flmCellsMaxWidth{0.96\linewidth}
\ExplSyntaxOn
\long\def\flm@cells@fit#1#2#3#4{%
  \setbox0=\hbox{\begin{tblr}{colspec={#1},#3}#4\end{tblr}}%
  \dim_compare:nNnTF {\wd0} < {\flmCellsMaxWidth}
    {\leavevmode\box0}
    {\begin{tblr}{width=\flmCellsMaxWidth,colspec={#2},#3}#4\end{tblr}}}
\long\def\flm@cells@iflr#1#2{\mode_if_horizontal:TF {\mode_if_inner:TF {#1} {#2}} {#2}}
\def\flm@cells@set#1{\ifcsname flmCellsKeys#1\endcsname
  \exp_args:Nnv \tblr_set_cell:nn {} {flmCellsKeys#1}\fi}
\ExplSyntaxOff
\ProvideDocumentCommand\flmCellsTable{+m +m +m +m}{%
  \flm@cells@iflr
    {\flm@cells@fit{#1}{#2}{#3}{#4}}%
    {\flmCellsBeginCenter\flm@cells@fit{#1}{#2}{#3}{#4}\flmCellsEndCenter}}
\ProvideDocumentCommand\flmCellsStyleKeys{m m}{%
  \expandafter\providecommand\csname flmCellsKeys#1\endcsname{#2}}
\ifcsname flmCellsSet\endcsname\else
  \NewTableCommand\flmCellsSet[1]{\flm@cells@set{#1}}
  \ProvideDocumentCommand\flmCellsSet{m}{}
\fi
\providecolor{flmCellsColorgreen}{RGB}{200,255,200}
\providecolor{flmCellsColorblue}{RGB}{200,220,255}
\providecolor{flmCellsColoryellow}{RGB}{255,255,200}
\providecolor{flmCellsColorred}{RGB}{255,200,200}
\flmCellsStyleKeys{green}{bg=flmCellsColorgreen}
\flmCellsStyleKeys{blue}{bg=flmCellsColorblue}
\flmCellsStyleKeys{yellow}{bg=flmCellsColoryellow}
\flmCellsStyleKeys{red}{bg=flmCellsColorred}
\endinput
